\begin{exercise}[Widerstand und Leistung]
\label{ex:resistance}
\given $U=\qty{12.0}{\volt}$, $I=\qty{0.150}{\ampere}$.
Es wird ein idealer ohmscher Widerstand bei konstanter Temperatur betrachtet.
\find Widerstand $R$ in \unit{\ohm} und elektrische Leistung $P$ in \unit{\watt}.
\solution
Mit den vorgegebenen Rechengrößen wird in \eqnref{eq:resistance} und
\eqnref{eq:power} gerechnet:
\begin{align}
  R &= \frac{U}{I}
     = \frac{\qty{12.0}{\volt}}{\qty{0.150}{\ampere}}
     = \qty{80.0}{\ohm}, \label{eq:resistance}\\
  P &= UI
     = \qty{12.0}{\volt}\cdot\qty{0.150}{\ampere}
     = \qty{1.80}{\watt}. \label{eq:power}
\end{align}
\result{R=\qty{80.0}{\ohm},\qquad P=\qty{1.80}{\watt}}
Die Werte sind frei gewählte Rechengrößen, keine Messdaten.
\end{exercise}

\begin{exercise}[Lineares Gleichungssystem]
\label{ex:system}
\given Das Gleichungssystem in \eqnref{eq:system}; $x$ und $y$ sind dimensionslos.
\find Die Unbekannten $x$ und $y$.
\begin{equation}
  \begin{aligned}
    2x+y &= 7,\\
    x-y  &= 2.
  \end{aligned}
  \label{eq:system}
\end{equation}
\solution
Durch Addition beider Gleichungen entfällt $y$:
\begin{align*}
  (2x+y)+(x-y) &= 7+2,\\
  3x &= 9 \quad\Longrightarrow\quad x=3.
\end{align*}
Durch Einsetzen in die zweite Gleichung wird $y$ bestimmt:
\begin{equation*}
  y=x-2=3-2=1.
\end{equation*}
\result{x=3,\qquad y=1}
Die Probe ergibt $2\cdot3+1=7$ und $3-1=2$.
\end{exercise}

\paragraph*{\NotesText{Formelzeichen}{Symbols}}
$U$: Spannung in \unit{\volt}; $I$: Stromstärke in \unit{\ampere};
$R$: Widerstand in \unit{\ohm}; $P$: Leistung in \unit{\watt};
$x$: erste Unbekannte, dimensionslos; $y$: zweite Unbekannte, dimensionslos.
